% Textbook Section 2.5, pp. 156--164.
\section[Factorizations]{Matrix Factorizations}
\sectioncard{2.5}{Matrix Factorizations}

\newcommand{\pivotentry}[1]{\colorbox{blue!15}{\(\mathstrut #1\)}}

\begin{frame}{Why Factor a Matrix?}
Matrix multiplication combines the effects of several transformations. A factorization analyzes a matrix by expressing it as a product of matrices with useful structure.
\pause
The LU factorization is motivated by repeated systems with the same coefficient matrix:
\[
A\vect x=\vect b_1,\qquad
A\vect x=\vect b_2,\qquad\ldots,\qquad
A\vect x=\vect b_p.
\]
\pause
Row reduction of the first system can also preprocess $A$. The resulting factors can then be reused for every later right-hand side.
\end{frame}

\begin{frame}{The LU Factorization}
Suppose row replacement operations reduce an $m\times n$ matrix $A$ to an echelon form $U$ without row interchanges. Then
\[
A=LU,
\]
where $L$ is an invertible $m\times m$ unit lower-triangular matrix.
\pause
To solve $A\vect x=\vect b$, write
\[
L(U\vect x)=\vect b.
\]
With $\vect y=U\vect x$, solve the two triangular systems
\[
\boxed{L\vect y=\vect b}
\qquad\text{followed by}\qquad
\boxed{U\vect x=\vect y}.
\]
\end{frame}

\begin{frame}[shrink=8]{Textbook Example 1: A Given LU Factorization}
\[
A=\begin{bmatrix}
3&-7&-2&2\\
-3&5&1&0\\
6&-4&0&-5\\
-9&5&-5&12
\end{bmatrix}
=
\underbrace{\begin{bmatrix}
1&0&0&0\\
-1&1&0&0\\
2&-5&1&0\\
-3&8&3&1
\end{bmatrix}}_{L}
\underbrace{\begin{bmatrix}
3&-7&-2&2\\
0&-2&-1&2\\
0&0&-1&1\\
0&0&0&-1
\end{bmatrix}}_{U}.
\]
Solve $A\vect x=\vect b$, where
\[
\vect b=\begin{bmatrix}-9\\5\\7\\11\end{bmatrix}.
\]
\pause
We first solve $L\vect y=\vect b$, and then solve $U\vect x=\vect y$.
\end{frame}

\begin{frame}{Example 1: Forward Substitution}
\small
\[
\left[\begin{array}{rrrr|r}
1&0&0&0&-9\\
-1&1&0&0&5\\
2&-5&1&0&7\\
-3&8&3&1&11
\end{array}\right]
\sim
\left[\begin{array}{rrrr|r}
1&0&0&0&-9\\
0&1&0&0&-4\\
0&0&1&0&5\\
0&0&0&1&1
\end{array}\right].
\]
\pause
Thus
\[
\vect y=\begin{bmatrix}-9\\-4\\5\\1\end{bmatrix}.
\]
Because $L$ is lower triangular, the entries of $\vect y$ are obtained from top to bottom.
\end{frame}

\begin{frame}{Example 1: Back Substitution}
\small
\[
\left[\begin{array}{rrrr|r}
3&-7&-2&2&-9\\
0&-2&-1&2&-4\\
0&0&-1&1&5\\
0&0&0&-1&1
\end{array}\right]
\sim
\left[\begin{array}{rrrr|r}
1&0&0&0&3\\
0&1&0&0&4\\
0&0&1&0&-6\\
0&0&0&1&-1
\end{array}\right].
\]
\pause
Hence
\[
\boxed{\vect x=\begin{bmatrix}3\\4\\-6\\-1\end{bmatrix}}.
\]
The entries are obtained from bottom to top because $U$ is upper triangular.
\end{frame}

\begin{frame}{Example 1: Verification and Work}
Substitution verifies
\[
A\begin{bmatrix}3\\4\\-6\\-1\end{bmatrix}
=\begin{bmatrix}-9\\5\\7\\11\end{bmatrix}.
\]
\pause
The textbook counts 28 arithmetic operations after $L$ and $U$ are known. Reducing the augmented matrix $[\,A\ \vect b\,]$ to $[\,I\ \vect x\,]$ takes 62 operations.
\medskip
\begin{block}{Reason for reuse}
For a new right-hand side, the same $L$ and $U$ are used again. Only the two triangular systems change.
\end{block}
\end{frame}

\begin{frame}{Textbook Example 2: Setup}
Find an LU factorization of
\[
A=\begin{bmatrix}
2&4&-1&5&-2\\
-4&-5&3&-8&1\\
2&-5&-4&1&8\\
-6&0&7&-3&1
\end{bmatrix}.
\]
Since $A$ has four rows, $L$ should be $4\times4$.
\end{frame}

\begin{frame}{Textbook Example 2: The First Column of $L$}
The first column of $L$ is the first column of $A$ divided by the top pivot entry:
\[
\frac{1}{2}\begin{bmatrix}2\\-4\\2\\-6\end{bmatrix}
=\begin{bmatrix}1\\-2\\1\\-3\end{bmatrix}.
\]
Thus
\[
L=\begin{bmatrix}
1&0&0&0\\
-2&1&0&0\\
1& &1&0\\
-3& & &1
\end{bmatrix},
\]
where the remaining entries will come from the later pivot columns.
\end{frame}

\begin{frame}[shrink=11]{Textbook Example 2: Row Reduction Produces $U$}
The highlighted entries determine the sequence of row operations from $A$ to $U$:
\[
A=
\begin{bmatrix}
\pivotentry{2}&4&-1&5&-2\\
\pivotentry{-4}&-5&3&-8&1\\
\pivotentry{2}&-5&-4&1&8\\
\pivotentry{-6}&0&7&-3&1
\end{bmatrix}
\sim A_1=
\begin{bmatrix}
2&4&-1&5&-2\\
0&\pivotentry{3}&1&2&-3\\
0&\pivotentry{-9}&-3&-4&10\\
0&\pivotentry{12}&4&12&-5
\end{bmatrix}
\]
\[
A_1\sim A_2=
\begin{bmatrix}
2&4&-1&5&-2\\
0&3&1&2&-3\\
0&0&0&\pivotentry{2}&1\\
0&0&0&\pivotentry{4}&7
\end{bmatrix}
\sim U=
\begin{bmatrix}
2&4&-1&5&-2\\
0&3&1&2&-3\\
0&0&0&2&1\\
0&0&0&0&\pivotentry{5}
\end{bmatrix}.
\]
\end{frame}

\begin{frame}{Textbook Example 2: Columns of $L$}
At each pivot column, divide the highlighted entries by the pivot and place the result into the corresponding column of $L$:
\[
\frac{1}{2}\begin{bmatrix}2\\-4\\2\\-6\end{bmatrix}
=\begin{bmatrix}1\\-2\\1\\-3\end{bmatrix},
\qquad
\frac{1}{3}\begin{bmatrix}3\\-9\\12\end{bmatrix}
=\begin{bmatrix}1\\-3\\4\end{bmatrix},
\]
\[
\frac{1}{2}\begin{bmatrix}2\\4\end{bmatrix}
=\begin{bmatrix}1\\2\end{bmatrix},
\qquad
\frac{1}{5}\begin{bmatrix}5\end{bmatrix}
=\begin{bmatrix}1\end{bmatrix}.
\]
The resulting columns begin at successive pivot rows.
\end{frame}

\begin{frame}{Textbook Example 2: The Factor $L$}
Placing the four resulting columns into a unit lower triangular matrix gives
\[
L=\begin{bmatrix}
1&0&0&0\\
-2&1&0&0\\
1&-3&1&0\\
-3&4&2&1
\end{bmatrix}.
\]
Together with the echelon form $U$ from the row reduction,
\[
\boxed{A=LU}.
\]
An easy calculation verifies that this $L$ and $U$ satisfy $LU=A$.
\end{frame}

\begin{frame}{Why the Factorization Works}
For the matrices constructed in Example 2, direct multiplication gives $LU=A$.
\pause
The same row operations reduce $A$ to $U$ and $L$ to $I$. If their elementary matrices are $E_1,\ldots,E_p$, then
\[
E_p\cdots E_1A=U,
\qquad E_p\cdots E_1L=I.
\]
\pause
Therefore $L=(E_p\cdots E_1)^{-1}$ and
\[
A=(E_p\cdots E_1)^{-1}U=LU.
\]
\end{frame}

\begin{frame}{Algorithm for an LU Factorization}
\begin{block}{Step 1}
Reduce $A$ to an echelon form $U$ using row replacement operations, if possible.
\end{block}
\pause
\begin{block}{Step 2}
Place entries in $L$ so that the same sequence of row operations reduces $L$ to $I$.
\end{block}
\medskip
The row operations reduce both matrices in parallel:
\[
A\longrightarrow U,
\qquad
L\longrightarrow I.
\]
Therefore $L^{-1}A=U$, so $A=LU$.
\end{frame}

\begin{frame}{Row Interchanges and Partial Pivoting}
The preceding construction assumes that row replacement operations are sufficient. In practical work, row interchanges are usually needed for numerical accuracy.
\pause
Partial pivoting chooses a pivot with the largest absolute value among the available entries in its column.
\medskip
With row interchanges, the textbook describes $L$ as \textbf{permuted lower triangular}: a rearrangement of its rows produces a unit lower-triangular matrix.
\medskip
The two-stage solution process still applies, but the triangular reduction follows the pivot order.
\end{frame}

\begin{frame}[shrink=4]{Numerical Notes}
For a dense $n\times n$ matrix with moderately large $n$:
\begin{enumerate}
  \item Computing an LU factorization takes about
  \[
  \frac{2n^3}{3}\text{ flops}.
  \]
  Computing $A^{-1}$ takes about $2n^3$ flops.
  \item Solving $L\vect y=\vect b$ and $U\vect x=\vect y$ takes about $2n^2$ flops.
  \item Multiplying $A^{-1}\vect b$ also takes about $2n^2$ flops, but errors from forming $A^{-1}$ may reduce accuracy.
  \item If $A$ is sparse, $L$ and $U$ may remain sparse while $A^{-1}$ is likely to be dense.
\end{enumerate}
\end{frame}

\begin{frame}{Textbook Example 3: A Ladder Network}
A series circuit and a shunt circuit have transfer matrices
\[
A_1=\begin{bmatrix}1&-R_1\\0&1\end{bmatrix},\qquad
A_2=\begin{bmatrix}1&0\\-1/R_2&1\end{bmatrix}.
\]
The series connection applies $A_1$ first and then $A_2$:
\[
A_2A_1=
\begin{bmatrix}
1&-R_1\\
-1/R_2&1+R_1/R_2
\end{bmatrix}.
\]
The order $A_2A_1$ records that the input passes through the series circuit first.
\end{frame}

\begin{frame}{Textbook Example 3: Design by Factorization}
The target transfer matrix is
\[
\begin{bmatrix}1&-8\\-.5&5\end{bmatrix}.
\]
\pause
We seek $R_1$ and $R_2$ such that
\[
\begin{bmatrix}
1&-R_1\\
-1/R_2&1+R_1/R_2
\end{bmatrix}
=\begin{bmatrix}1&-8\\-.5&5\end{bmatrix}.
\]
\pause
Comparing entries gives
\[
R_1=8\ \Omega,\qquad \frac1{R_2}=.5\quad\Longrightarrow\quad R_2=2\ \Omega.
\]
The matrix factorization identifies the physical components required by the design.
\end{frame}
